\section{AI 是放大器：从学习反馈到认知卸载}

赵晨阳说 AI 是放大器：它放大人的主观诉求和主观能动性。这个判断有两面。对主动学习者来说，AI 是天大的福报。比如高中时想理解洛必达法则，传统教学可能只说“超纲不能用”，却不认真解释何时能用、为何能用；AI 可以快速回应求知欲，给出即时反馈，帮助人深入探索。

但尚晏仪提醒，要警惕认知卸载。如果一个人还没有想清楚自己要学什么、做什么，就把问题抛给 AI，让 AI 给多轮交付，长期会丧失主动思考能力。AI 反馈非常及时、甚至会哄人，图片和视频反馈又更刺激，很容易迅速建立多巴胺通路。她提到自己玩 AI 视频生成时感到这种反馈通路建立得非常快，因此觉得可怕。

\begin{warningbox}{AI 作为放大器，不自动放大好东西}
AI 会放大主观能动性，也会放大懒惰、冲动和短反馈成瘾。对于已经有清晰目标的人，它是学习加速器；对于没有目标和判断力的人，它可能成为认知卸载工具。
\end{warningbox}

Sora 下架和视频生成商业模式的讨论，延续了这一点。赵晨阳认为 Sora 技术并不差，问题可能在于视频推理太贵，而 OpenAI 自己没有视频生态能把生成内容融入 revenue 循环。快手或字节这类本来就有视频生态的平台，生成视频可以进入自己的分发和变现闭环；OpenAI 做一个独立视频产品，则会面对成本与收入不匹配。

这一段其实把“AI 是放大器”从个人学习扩展到商业生态。AI 视频会放大用户娱乐欲望，也会放大平台分发能力；但若没有生态闭环，技术强也未必是好生意。相同技术在不同组织结构中会产生完全不同的商业结果。

\begin{knowledgebox}{放大器逻辑的三层}
\begin{tabularx}{\textwidth}{>{\raggedright\arraybackslash}p{0.2\textwidth} Y}
\toprule
层次 & 放大什么 \\
\midrule
个人学习 & 放大求知欲、反馈速度和自我训练能力，也可能放大认知卸载。 \\
代码生产 & 放大执行速度，也放大低质量 PR 和 review 压力。 \\
平台生态 & 放大内容生成和分发能力，但前提是成本能进入商业闭环。 \\
\bottomrule
\end{tabularx}
\end{knowledgebox}

\subsection{本章小结}

AI 不是单纯“更聪明的工具”，而是放大器。它会放大主动学习者的探索，也会放大不思考者的依赖；会放大代码生产，也会放大责任缺口；会放大视频生成，也会放大平台成本和内容刺激。关键不是能不能用 AI，而是使用者有没有目标、审美和边界。

